\section{Introduction}
\label{sec:intro}

Picture an SSH session at 03:14 on a Tuesday. An attacker runs
\texttt{nmap}, sprays credentials, lands inside a Cowrie~2.5 container, and
types \texttt{uname -a}. Cowrie logs the keystrokes, timestamps them, and
returns to waiting. For two decades this has been the state of the art: a
high-fidelity recorder that never reasons about what it records. A recent
survey of production deception frameworks found that no widely deployed
honeypot grounds its responses in real-time threat intelligence or
differentiates its behavior across attacker classes~\cite{lopez2024}. The
recorder watches; it does not think.

That silence is expensive. The mean time to identify a breach is measured in
months~\cite{ibm2024}, while adversary breakout time has compressed to
minutes~\cite{crowdstrike2026}. A honeypot sits directly in the path of the
attacker, yet converts almost none of that proximity into understanding. It
cannot tell a credential-stuffing bot from a hands-on-keyboard operator, it
serves the same canned banner to both, and it produces none of the
tactics-and-techniques intelligence that a security operations center would
need to act on the encounter.

Large language models (LLMs) promise to close this gap. Prompted with a
session transcript and a library of threat intelligence, a modern LLM can
impersonate a vulnerable Linux host, improvise plausible file systems, and
steer the conversation toward the attacker's next objective. Recent systems
have begun to exploit this, deploying LLM agents as honeypots that chat an
attacker through a fictitious machine~\cite{reworr2024}. The promise,
however, arrives with a problem the deception literature has largely left
unexamined: \emph{the model can be wrong in ways a script cannot}. A
honeypot's product is a lie, and the lie must satisfy three properties at
once. It must be \emph{consistent}: the kernel version printed on turn one
cannot contradict the one an attacker reads on turn five, and the database
banner cannot disagree with the port the same host just advertised. It
must be \emph{safe to serve}: a model that complies too well will invent
credential material with no way to distinguish what it fabricated from
what it leaked, poisoning every downstream analytical product. And it
must be \emph{available}: a generator is a rented service, and ours was.
Under safety-alignment pressure from one commercial provider, 86.7\% of
generated turns exceeded the inference budget, and response fidelity
depended on how much quota the account happened to have. Unverified
generation can forfeit all three properties in the same session.

A natural question arises: why not fix this inside the model? Careful
prompting, refusal training, and constrained decoding all shape
probabilities, but none can certify a particular output. Alignment
training produces models whose refusal behavior shifts with phrasing and
pressure~\cite{ganguli2022,bai2022}, which is precisely the failure our
quota cliff exposed; and a deceived honeypot is itself an LLM application
with attacker-controlled context, so injection techniques that hijack
system prompts~\cite{greshake2023} apply with the attacker already in the
room. A probability, however small, is not a promise a security tool
can serve on. The assurance has to live outside the model, in the
pipeline around it.

The classical system-side fixes fail differently. Pre-scripting every
answer regresses to the finite automaton of low-interaction deception,
discarding the adaptivity that motivated generation in the first place;
real-time human review of every response cannot keep pace with an
attacker who reads faster than an analyst can vet. What survives is an
asymmetry between the two halves of the problem. Generating a plausible
answer to an arbitrary question is open-ended: its output space cannot
be enumerated, so no amount of testing certifies the next draw.
Checking a draft against a small, closed set of claims is the opposite:
a bounded computation whose verdict is final. A system can exploit this
by splitting the roles, letting the generator propose and a verifier
dispose against a fiction the system itself fully controls.

This paper takes the position that the right question is not
\emph{can an LLM impersonate a host} but \emph{how much intelligence can
a lie safely spend, and where must the spending stop}. Deception systems
build intelligence in rungs -- templates, classifiers, retrieval, free
generation -- and each rung has a trust price that practice pays
implicitly. We make the pricing explicit. Our answer is a verified
deception loop in which
each mechanism exists to buy one of the three properties, placed on a
single axis of certifiability (\S\ref{sec:ladder}). A fiction graph
binds every host fact to one consistent set of triples, and a verifier
checks every draft against that graph before anything reaches the wire:
consistency, enforced at serving time, not asked for at prompt time.
Baited content is drawn only from canary-registered secrets, so that any
disclosure is a measured event rather than a leak: safety by
construction. A deterministic fallback serves human-vetted transcripts
whenever generation or verification fails, so the system degrades to the
last known good answer rather than an unverified improvisation:
availability under provider failure. Beneath these sit a behavioral
profiler and a hybrid retrieval engine that ground a per-attacker
deception plan in curated CTI, giving the lie its targeting, and a
human-in-the-loop gate holds back any content class marked as requiring
approval. The loop turns generation from an act of trust into a measured
pipeline with an audit trail, and the evaluation reads the whole
apparatus as an escalation study: each rung of intelligence is measured
for what it adds over the rung beneath it, including the rungs that fold
back down when a certifiable mechanism reproduces their effect.

We evaluate whether each element of this loop contributes measurably,
using
artifact-backed benchmarks that we release alongside the paper. Four
questions structure the study:
\begin{itemize}
\item \textbf{RQ1:} Does retrieval grounding measurably change what the
honeypot detects and serves? A paired, per-scenario ablation across 90
matched sessions finds that grounding lifts technique detection by
$+0.317$ (95\% CI $[+0.162, +0.480]$) and actor attribution by $+0.140$
($[+0.083, +0.197]$), while leaving persona realism untouched, a null
result the architecture predicts and one we treat as evidence the
ablation discriminates.
\item \textbf{RQ2:} Can an LLM-driven honeypot remain internally consistent
and safe under sustained adversarial interaction? Across 45 targeted
re-queries the fiction graph serves zero contradictions with byte-identical
determinism, and across a 218-query live campaign the verification layer
records no unverified response served and no false-positive CTI alert.
\item \textbf{RQ3:} Is the resulting deception strong enough to hold a
skilled observer? A three-evaluator consensus study, each evaluator an
independently prompted LLM reviewer persona distinct from the generator,
scores the full session transcripts at $3.86/5$ (exact agreement
$0.750$, a within-model consistency measure -- all three personas share
one judge model), with all eight scenarios above the acceptance bar and
a known judge bias corrected by canary-registered secrets.
\item \textbf{RQ4:} Does each additional rung of intelligence justify its
trust price? The escalation analysis adds each mechanism
to the same
loop and prices its marginal effect: routing accuracy rises from 0.8875
to 0.9550; grounding contributes the $+0.154$ composite above; and at
the top rung, kind-aware adversarial classification corrects 12/13
real-world assessments and raises attacker bait exposure by 125\% in
closed-loop replay, yet adds no measurable dwell time and is fully
measurably subsumed by a deterministic decision floor at the benchmark's
resolution -- the fold-down result that turns the ladder from a build
history into a design rule.
\end{itemize}

\noindent Our contributions are:
\begin{itemize}
\item \textbf{A verified deception architecture} (\S\ref{sec:design}) that
couples RAG-grounded generation to a fiction graph and a serving-time
verification gate, with deterministic fallback, so that every served
response is either verified against the deception state or produced from
vetted static content.
\item \textbf{A measurable grounding effect} (\S\ref{sec:evaluation}):
paired-difference bootstrap intervals over identical scenario sets show
$+0.154$ overall composite improvement from retrieval grounding
($[+0.092, +0.221]$), with 35 TTPs detected against 0 without grounding.
\item \textbf{A consistency and safety benchmark} for LLM honeypots
(\S\ref{sec:evaluation}): contradiction-free serving under adversarial
re-query, next-technique steering accuracy (top-3 $0.9$), and an ablation
battery (A1--A5) in which each removed mechanism costs the system the
property it was added to provide.
\item \textbf{Run-to-run stability evidence} rarely reported for LLM
systems: two independent 218-query campaigns reproduce the composite to
$+0.005$ with $0.991$ per-row pass agreement, the deterministic
strata reproduce exactly, and the comparison survives unchanged when
recomputed on the judge-audited rows.
\item \textbf{An open, reproducible artifact} on commodity hardware
(2~vCPU, ${\sim}4$~GB RAM), with every number in this paper traceable to a
persisted result file.
\item \textbf{An escalation methodology for AI in deception}
(\S\ref{sec:ladder}, \S\ref{sec:jeveval}): mechanisms placed on a
certifiability axis, each rung priced by measured marginal effect, and
the descent direction exercised on live machinery -- when a certifiable
mechanism reproduces an intelligent one's benefit, the intelligent one
comes out.
\item \textbf{An audited measurement pipeline} (\S\ref{sec:evaluation}):
the automated judge itself is audited offline against the transcripts it
scored, a full-corpus census locates the stratum where the fold-down
comparison is informative, and the audit-plus-stability chain reseeds
automatically whenever a new campaign run lands.
\item \textbf{A scorer-validity audit and declared measurement channels}
(\S\ref{sec:scorer-audit}): a double-annotated, blinded audit of the
automated rubric against a hash-pinned \attck{} taxonomy, which caught a
realized scorer-validity failure (echo-crediting on fallback-heavy
serves) and motivated a correction that moves machine-vs-human agreement
from $\kappa=-0.187$ to $+0.349$; served-text and analysis metrics are
reported as declared, separate channels, with an external-baseline ladder
(vanilla Cowrie, stateful-LLM, bare-LLM, grounded-without-verifier, full
verified loop) on one workload.
\item \textbf{Cross-session consistency and a timing-fingerprint self-assessment}
(\S\ref{sec:rq2}, \S\ref{sec:discussion}): a host-level fiction store
catches cross-session drift that per-session verification cannot see by
construction (4/4 vs.\ 1/4), at zero paired cost on benign workloads; and
the paper measures its own timing side channel (AUC $0.983$ unmitigated)
with a mitigation that reaches the drift floor.
\end{itemize}

We are careful about what this paper does not claim. The deployment is
tailnet-only, so it observes no ambient internet background radiation; the
evaluation judge shares a provider with, but is a different model from, the
generator; and the system's availability on free-tier inference is a
constraint any practitioner should weigh. These limits are quantified
where they arise (\S\ref{sec:discussion}).

The remainder of the paper is organized as follows.
Section~\ref{sec:background} fixes vocabulary and the threat model.
Section~\ref{sec:design} presents the verified deception loop.
Section~\ref{sec:implementation} gives implementation details.
Section~\ref{sec:evaluation} reports the benchmark battery, the ablations,
the consensus study, and the stability runs; Appendix~\ref{app:artifacts}
indexes the released artifacts behind every number.
Section~\ref{sec:discussion} discusses limitations and deployment
guidance, Section~\ref{sec:related} positions the work in the deception
and LLM-security literature, and Section~\ref{sec:conclusion} concludes.
